\chapter{Inducción y algoritmos recursivos}
\label{md:induccion}

\section{Inducción ordinaria y fuerte}

Para demostrar una propiedad $P(n)$ para todos los enteros $n\ge n_0$,
el \textbf{principio de inducción} exige:
\begin{enumerate}
 \item \textbf{Caso inicial:} comprobar $P(n_0)$.
 \item \textbf{Paso inductivo:} para un $n\ge n_0$ arbitrario, suponer
       $P(n)$ y deducir $P(n+1)$.
\end{enumerate}
La hipótesis inductiva es una suposición dentro del paso, no la afirmación
de que el resultado ya está demostrado. El caso inicial y el paso
permiten alcanzar sucesivamente todos los índices del intervalo.

La \textbf{inducción fuerte} permite suponer $P(n_0),\ldots,P(n)$
para deducir $P(n+1)$. Es útil cuando el siguiente caso depende de varios
anteriores, o de un índice menor que no tiene por qué ser $n$.
Ambas formas se justifican por buena ordenación: si existiera un
contraejemplo, el menor contradiría el caso inicial o el paso inductivo.

Si un paso utiliza dos o más términos previos, se comprueban todos
los casos iniciales necesarios. Un paso que solo llega de $n$ a $n+2$
no cubre por sí solo las dos paridades.

\section{Definiciones y procedimientos recursivos}

Una definición recursiva especifica casos iniciales y una regla para
obtener nuevos valores a partir de otros ya definidos.
La regla debe alcanzar los datos de partida tras un número finito de
reducciones para que el cálculo termine.

El factorial ilustra este esquema:
\[
 0!=1,\qquad n!=n(n-1)!\quad(n\ge1).
\]
Para calcular $n!$ se reduce el argumento hasta $0$ y después se
reconstruyen los productos. La inducción comprueba que el procedimiento
produce el producto $1\cdot2\cdots n$.

La corrección de un algoritmo recursivo requiere comprobar el resultado
en los casos base y que la reducción, suponiendo correctos los subcasos,
produce el resultado requerido. La terminación se prueba aparte mediante
una cantidad entera no negativa que decrezca en cada llamada.
Euclides (sección~\ref{md:euclides}) es un ejemplo: el resto disminuye.
Las relaciones de recurrencia y su resolución se estudian en el
capítulo~\ref{md:recurrencias}.
